%! AP Statistics — Unit 2, Lesson 2.7 — Group Activity
\documentclass[10pt]{article}
\usepackage{apstats-article}
\usepackage{apstats-boxes}

\begin{document}

\pageheader{Unit 2, Lesson 2.7}{Group Activity}
\namepartnerperiod

\vspace{0.18in}

% ── SCENARIO ─────────────────────────────────────────────────────────────────
\begin{scenariobox}[Directions]{sky}
\small
Each tier uses a real-world regression context. Compute residuals, construct or interpret a
residual plot, and write a complete AP-style conclusion about model appropriateness. Write
all sentences in context using the variable names.
\end{scenariobox}

\vspace{0.12in}

% ── TIER R ────────────────────────────────────────────────────────────────────
\begin{tcolorbox}[colback=white, colframe=black!40, boxrule=0.8pt, arc=2mm,
    left=4mm, right=4mm, top=3mm, bottom=3mm,
    title={\textbf{Tier R --- Remediate}}, fonttitle=\bfseries, halign title=center]
\small

\textbf{Context: Age and Blood Pressure}

A health clinic collected data on patient age ($x$, years) and systolic blood pressure ($y$,
mmHg). Technology produced the regression equation:
\[
  \widehat{\text{BP}} = 80 + 0.5(\text{age})
\]
The table below shows age, blood pressure, and predicted blood pressure for five patients.

\vspace{8pt}
\renewcommand{\arraystretch}{1.9}
\begin{tabularx}{\linewidth}{@{} >{\centering\arraybackslash}p{2.4cm}
                                    >{\centering\arraybackslash}p{2.4cm}
                                    >{\centering\arraybackslash}p{2.8cm}
                                    X @{}}
\rowcolor{sky}
\textbf{Age ($x$)} & \textbf{BP ($y$)} & \textbf{Predicted ($\hat{y}$)} &
\textbf{Residual ($e = y - \hat{y}$)} \\
\midrule
20 & 90 & 90 & \blank{2.5cm} \\
30 & 96 & 95 & \blank{2.5cm} \\
40 & 99 & 100 & \blank{2.5cm} \\
50 & 106 & 105 & \blank{2.5cm} \\
60 & 109 & 110 & \blank{2.5cm} \\
\midrule
\multicolumn{3}{r}{\textbf{Sum of residuals}} & \blank{2.5cm} \\
\end{tabularx}

\vspace{10pt}

\begin{enumerate}[leftmargin=1.6em, itemsep=0.18in]

  \item Compute and record the residual for each patient. Verify that the residuals sum to 0.\\[6pt]
        (Check: sum $=$ \blank{2cm})

  \item Interpret the residual for the 30-year-old patient in context.\\[8pt]
        \writeline\\[3pt]\writeline\\[3pt]\writeline

  \item A residual plot for these data shows random scatter around $e = 0$ with no obvious
        pattern. Write a complete conclusion about whether a linear model is appropriate for
        predicting blood pressure from age.\\[8pt]
        \writeline\\[3pt]\writeline\\[3pt]\writeline

\end{enumerate}
\end{tcolorbox}

\vspace{0.12in}

% ── TIER A ────────────────────────────────────────────────────────────────────
\begin{tcolorbox}[colback=white, colframe=black!40, boxrule=0.8pt, arc=2mm,
    left=4mm, right=4mm, top=3mm, bottom=3mm,
    title={\textbf{Tier A --- Approaching Proficiency}}, fonttitle=\bfseries, halign title=center]
\small

\textbf{Context: Hours of TV and GPA}

A researcher recorded daily TV-watching hours ($x$) and GPA ($y$) for five students.
The data ranged from \textbf{1 to 5 hours} of TV per day. Technology output:
\[
  \widehat{\text{GPA}} = 4.05 - 0.35(\text{TV hours})
\]

\vspace{8pt}
\renewcommand{\arraystretch}{1.9}
\begin{tabularx}{\linewidth}{@{} >{\centering\arraybackslash}p{2.6cm}
                                    >{\centering\arraybackslash}p{2.0cm}
                                    >{\centering\arraybackslash}p{2.8cm}
                                    X @{}}
\rowcolor{sky}
\textbf{TV hours ($x$)} & \textbf{GPA ($y$)} & \textbf{Predicted ($\hat{y}$)} &
\textbf{Residual ($e = y - \hat{y}$)} \\
\midrule
1 & 3.90 & \blank{2cm} & \blank{2.5cm} \\
2 & 3.00 & \blank{2cm} & \blank{2.5cm} \\
3 & 3.20 & \blank{2cm} & \blank{2.5cm} \\
4 & 2.50 & \blank{2cm} & \blank{2.5cm} \\
5 & 2.40 & \blank{2cm} & \blank{2.5cm} \\
\midrule
\multicolumn{3}{r}{\textbf{Sum}} & \blank{2.5cm} \\
\end{tabularx}

\vspace{10pt}

\begin{enumerate}[leftmargin=1.6em, itemsep=0.18in]

  \item Compute $\hat{y}$ and $e$ for each student. Verify $\sum e = 0$ (allow rounding to 0.01).

  \item Plot the residuals on the axes below. Draw a dashed reference line at $e = 0$.

        \vspace{0.1in}
        \begin{center}
        \begin{tikzpicture}
          \draw[->] (0,0) -- (7,0) node[right]{\small TV hours ($x$)};
          \draw[->] (0,-1.6) -- (0,1.6) node[above]{\small residual};
          \draw[dashed,gray] (0,0) -- (7,0);
          \foreach \x/\lbl in {1.2/1, 2.4/2, 3.6/3, 4.8/4, 6.0/5}
            \draw (\x,0.06)--(\x,-0.06) node[below]{\tiny\lbl};
          \foreach \y/\lbl in {-1.2/{$-0.4$}, -0.6/{$-0.2$}, 0.6/{$0.2$}, 1.2/{$0.4$}}
            \draw (0.06,\y)--(-0.06,\y) node[left]{\tiny\lbl};
        \end{tikzpicture}
        \end{center}

  \item Describe the pattern in the residual plot. Write a complete AP-style conclusion about
        whether a linear model is appropriate for predicting GPA from TV-watching hours.\\[8pt]
        \writeline\\[3pt]\writeline\\[3pt]\writeline

\end{enumerate}
\end{tcolorbox}

\vspace{0.12in}

% ── TIER E ────────────────────────────────────────────────────────────────────
\begin{tcolorbox}[colback=white, colframe=black!40, boxrule=0.8pt, arc=2mm,
    left=4mm, right=4mm, top=3mm, bottom=3mm,
    title={\textbf{Tier E --- Extension}}, fonttitle=\bfseries, halign title=center]
\small

\textbf{Context: Fertilizer and Plant Growth}

A botanist applied different amounts of fertilizer ($x$, grams per week) to seedlings and
measured their height ($y$, cm) after four weeks. Two different regression models were fit.
Their residual plots are described below.

\vspace{8pt}

\textbf{Model 1 --- Linear:} The residual plot shows a clear U-shaped (curved) pattern.
Residuals are large and positive at low and high fertilizer values, and negative in the middle.

\textbf{Model 2 --- Quadratic:} The residual plot shows random scatter around zero with no
obvious pattern.

\vspace{10pt}

\begin{enumerate}[leftmargin=1.6em, itemsep=0.18in]

  \item Write a complete AP-style statement about whether the linear model is an appropriate
        model for predicting plant height from fertilizer amount. Include the pattern and conclusion.\\[8pt]
        \writeline\\[3pt]\writeline\\[3pt]\writeline

  \item Write a complete AP-style statement about the quadratic model.\\[8pt]
        \writeline\\[3pt]\writeline\\[3pt]\writeline

  \item Which model would you recommend and why? Reference both residual plots in your answer.\\[8pt]
        \writeline\\[3pt]\writeline\\[3pt]\writeline

  \item The botanist's data shows the largest absolute residual at a fertilizer amount of
        $x = 12$ grams: the actual height was 32 cm and $\hat{y} = 26$ cm. Calculate and
        interpret this residual in context.\\[8pt]
        \writeline\\[3pt]\writeline

\end{enumerate}
\end{tcolorbox}

\end{document}
